\section{Result 1: fault-domain concentration creates an attribution boundary}
\label{sec:attribution}

This result uses only the Gray Paper verdict composition rules plus the fault-domain model in \Cref{sec:model}. It does not use the ELVES adaptive-crash branching approximation.

\subsection{Exact finite-population verdict conditions}

Suppose enough judgments eventually exist to construct a verdict. Let \(P\) be the number of validators that would issue a positive judgment for the invalid report and \(N-P\) the number issuing a negative judgment.

A bad verdict can be constructed if and only if at least \(K\) negative judgments are available:
\[
\boxed{
N-P\ge K.
}
\]
A good verdict can be constructed if and only if
\[
\boxed{
P\ge K.
}
\]
A wonky verdict can be constructed if and only if at least \(W\) positive and \(K-W\) negative judgments are available:
\[
\boxed{
P\ge W
\quad\text{and}\quad
N-P\ge K-W.
}
\]
These are selection conditions: a block author can choose the required \(K\) signed judgments from the available pool.

For JAM's current \(N=1023\),
\[
K=683,
\qquad
W=341,
\qquad
K-W=342.
\]
Therefore:
\begin{itemize}
\item a bad verdict is possible for \(P\le340\);
\item a wonky verdict is possible for \(341\le P\le681\);
\item a good verdict is possible for \(P\ge683\).
\end{itemize}

Under the literal current rule, \(P=682\) is a one-vote edge case in which none of the three allowed positive counts can be selected from a 683-judgment verdict: there are only 341 negatives, too few for the wonky composition, but only 682 positives, too few for a good verdict. We treat this as a specification question for the JAM authors rather than silently smoothing it away.

\subsection{Robust bad-verdict attribution}

Under correlated honest failure, honest buggy validators contribute positive judgments, honest independent validators contribute negative judgments, and adversarial validators may choose either sign.

To prevent a bad verdict, the adversary's best choice is to vote positive. The maximum positive population is then
\[
P_{\max}=A+B_x,
\]
and the remaining negative population is \(C_x\).

Hence a bad verdict remains available regardless of adversarial voting if and only if
\[
\boxed{
C_x\ge K.
}
\]
Equivalently,
\[
A+B_x\le N-K.
\]

Writing
\[
p_x=\frac{A+B_x}{N}=\gamma+(1-\gamma)f_x,
\]
the large-\(N\) form is
\[
\boxed{p_x<\frac13.}
\]
Solving for \(f_x\) gives
\[
\boxed{
 f_x<\frac{1/3-\gamma}{1-\gamma},
}
\]
when \(\gamma<1/3\).

At \(\gamma=0.2\), the asymptotic limit is
\[
f_x<\frac{1/3-0.2}{0.8}=\frac16\approx0.167.
\]
Thus a shared fault affecting roughly one-sixth of the honest population can consume the entire remaining one-third attribution budget when the adversarial share is 20\%.

The interpretation is narrow but important. Below this boundary, independently correct honest validators alone supply enough negative judgments for a bad verdict. Above it, the adversary can prevent a bad verdict simply by joining the buggy positive bloc.

\subsection{The wonky region}

Suppose
\[
A+B_x\ge W
\]
but independently correct negative validators are still numerous enough for the wonky composition:
\[
C_x\ge K-W.
\]
Then an adversary voting positive can make a wonky verdict constructible while preventing a bad verdict.

Under the current Gray Paper, a wonky verdict is non-positive and therefore stops the report from proceeding, but the culprit/fault rules do not create ordinary offender records for wonky reports. Thus the system can remain \emph{safety preserving at the report level} while losing the ordinary attribution path.

This is the first important separation:
\[
\boxed{
\text{report rejection}
\neq
\text{offender attribution}.
}
\]

It is deliberately not stated as “the attacker pays nothing.” The staking layer may define additional consequences for wonky records. The Gray Paper does not specify such a rule, so whether wonky is economically cost-free is an explicit question for the JAM/Polkadot staking design.

\subsection{Strategic choice inside the intermediate region}

There is a second asymmetry. If the honest buggy population by itself is still small enough that
\[
B_x\le N-K,
\]
the adversary can instead vote negative. Then at least \(K\) negative judgments may be available, making a bad verdict constructible. In that branch the honest buggy validators' positive judgments contradict the bad verdict and their guarantor/judgment signatures may become offender evidence, whereas adversarial validators that voted negative do not contradict it.

Thus in part of the intermediate region the adversary can choose between two qualitatively different protocol outcomes:
\begin{enumerate}[label=(\alph*)]
\item vote positive, preventing a bad verdict and steering toward wonky; or
\item vote negative, permitting a bad verdict whose fault attribution can fall on honest buggy validators.
\end{enumerate}

This does not prove an attacker can always avoid all economic loss, because entering the pipeline may itself require attacker-controlled signatures or other costs. It shows that the verdict layer does not by itself force the collateral cost to fall on the party exploiting the shared fault.

\subsection{False-good eligibility}

If
\[
A+B_x\ge K,
\]
then a good verdict becomes constructible when the adversarial and buggy-positive blocs align. Asymptotically,
\[
\boxed{p_x>\frac23.}
\]
For \(\gamma=0.2\), this occurs at approximately
\[
f_x>\frac{2/3-0.2}{0.8}\approx0.583.
\]
At that point the dispute layer itself can accept the invalid report as good under the fault model, and correct negative judges fall on the dissenting side of the protocol verdict.

\subsection{Threshold ordering}

For JAM's \(F=2\), \Cref{sec:escalation} will show that the ELVES corrective branching process loses supercriticality at
\[
p_x=1-\frac1F=\frac12.
\]
The three large-population boundaries are therefore ordered as
\[
\boxed{
\frac13
<
\frac12
<
\frac23.
}
\]

\begin{figure}[htbp]
\centering
\begin{tikzpicture}[x=12cm,y=1cm,>=Latex,font=\small]
\draw[very thick] (0,0) -- (1,0);
\foreach \x in {0.333,0.5,0.667} {\draw[thick] (\x,0.12) -- (\x,-0.12);}
\node[align=center,font=\scriptsize] at (0.333,0.52) {attribution\\$p_x=1/3$};
\node[align=center,font=\scriptsize] at (0.5,0.82) {escalation\\$p_x=1/2$};
\node[align=center,font=\scriptsize] at (0.667,0.52) {false-good\\$p_x=2/3$};

\node[align=center,text width=1.7cm,font=\tiny] at (0.166,-0.55)
  {bad attribution robust};
\node[align=center,text width=1.8cm,font=\tiny] at (0.416,-0.55)
  {wonky can replace bad; audit supercritical};
\node[align=center,text width=1.8cm,font=\tiny] at (0.583,-0.55)
  {wonky region; audit subcritical};
\node[align=center,text width=1.7cm,font=\tiny] at (0.833,-0.55)
  {good verdict eligible};
\node[below=1.15cm] at (0.5,0)
  {$p_x=\gamma+(1-\gamma)f_x$: maximum wrong-positive share for fault $x$};
\end{tikzpicture}

\caption{Large-population fault-domain regimes for JAM's current \(F=2\). The first boundary concerns robust availability of a bad verdict and the ordinary offender-attribution path. The second is the ELVES-extension branching threshold. The third is eligibility for a good verdict under aligned adversarial and buggy-positive judgments. Exact finite-\(N\) rules use \(K=\floor{2N/3}+1\) and \(W=\floor{N/3}\).}
\label{fig:thresholds}
\end{figure}

At \(\gamma=0.2\), expressed as buggy share among otherwise-honest validators, these correspond approximately to 16.7\%, 37.5\%, and 58.3\% respectively.

\subsection{Offline validators tighten the attribution condition}

The verdict rule depends on available signed judgments, not only nominal validator identities. Let \(O\) validators be offline and unable to judge, disjoint from the adversarial and correlated-fault classes for this extension. Let \(f_x\) now denote the buggy share among the remaining non-adversarial online validators. The independently correct online share is
\[
\frac{C_x}{N}=(1-\gamma-o)(1-f_x),
\qquad o=\frac{O}{N}.
\]
A bad verdict robust to adversarial positive voting requires
\[
\boxed{
(1-\gamma-o)(1-f_x)\ge\frac{K}{N},
}
\]
which approaches
\[
(1-\gamma-o)(1-f_x)>\frac23.
\]
Offline capacity therefore consumes the same finite pool of negative judgments needed for attribution. This extension should not be conflated with ELVES no-shows: a validator that never announces an audit is not the same event as an announced auditor failing to return a judgment.
